\section{Signal Encoding: DenoMAE \& CDRF}

\begin{frame}{Background \& Related Work}
    \begin{itemize}
        \item Traditional AMC pipelines: supervised CNNs/ViTs on constellations, little reuse across datasets.
        \item CDRF benchmark showed YOLO/ResNet (pretrained on COCO/ImageNet) transfer to RF spectrograms.
        \item Need self-supervised encoder that internalizes noise statistics and scales with unlabeled RF data.
    \end{itemize}
\end{frame}

\begin{frame}{Research Question 1}
    \centering
    \Large
    \textit{Can vision-based, self-supervised learning effectively transfer world knowledge to the wireless domain?}
\end{frame}

\begin{frame}{CDRF as Empirical Foundation}
    \begin{itemize}
        \item 39 drone/RC classes, 500+ GB raw IQ, controlled SNR and interferer augmentations.
        \item YOLOv11n achieves mAP@0.5 up to 0.986 on clean spectrograms, validating cross-modal transfer.
        \item Provides realistic noise/interference motifs that motivate noise-aware pretraining.
    \end{itemize}
\end{frame}

\begin{frame}{DenoMAE Overview}
    \begin{itemize}
        \item Multimodal masked autoencoder tailored to RF constellations and IQ-derived images.
        \item Treats \textbf{noise} as an explicit modality to disentangle channel effects.
        \item Encoder is reused; decoders discarded post-pretraining for efficient deployment.
    \end{itemize}
\end{frame}

\begin{frame}{Architecture Pipeline}
    \begin{enumerate}
        \item Patchify $224\times224$ inputs (clean/noisy constellations, clean/noisy signals, noise maps).
        \item Mask 75\% of patches; process visibles via 12-layer ViT encoder ($d_{model}=768$).
        \item Project shared latent into modality-specific 4-layer decoders to reconstruct masked content.
    \end{enumerate}
\end{frame}

\begin{frame}{Pretraining Strategy}
    \begin{itemize}
        \item 10k unlabeled samples covering 10 modulation formats (BPSK, QPSK, 8PSK, QAM, GFSK, etc.).
        \item Masking enforces global reasoning, lowering compute yet boosting denoising skill.
        \item Optimized with AdamW, batch 64, learning rate $10^{-4}$, 100 epochs.
    \end{itemize}
\end{frame}

\begin{frame}{Data Pipeline}
    \begin{itemize}
        \item IQ samples reshaped to $32\times32$, interpolated, replicated to RGB for ViT compatibility.
        \item Noise modality generated via exponential decay model capturing neighborhood energy.
        \item Aligns with CDRF’s spectrogram evidence, creating an end-to-end cross-modal story.
    \end{itemize}
\end{frame}

\begin{frame}{Results: Robustness}
    \begin{itemize}
        \item +22.1\% accuracy at -10 dB SNR relative to non-pretrained ViT baseline.
        \item Maintains 77.5–84.3\% accuracy from -10 to +10 dB using only 1k labeled samples.
        \item Extrapolates down to -20 dB despite training $\geq$ -10 dB (see Figure~\ref{fig:extrap}).
    \end{itemize}
\end{frame}

\begin{frame}{Results: Efficiency}
    \begin{itemize}
        \item Needs 10× fewer unlabeled samples than AlexNet AMC and 3× fewer labeled samples at fine-tune.
        \item Ablations show noise modality delivers ~2\% extra accuracy, highlighting benefit of treating noise as data.
        \item Outperforms CNN-AMC, NMformer, DL-GRF on accuracy while using smaller datasets.
    \end{itemize}
\end{frame}

\begin{frame}{Ablation Highlights}
    \begin{itemize}
        \item Modalities removed sequentially: still 81.3\% with constellations only, 83.5\% with all five.
        \item Clean vs noisy constellation pairing adds +0.6\%; adding time-domain signal adds +1.3\%.
        \item Explicit noise branch regularizes encoder during pretraining without affecting inference cost.
    \end{itemize}
\end{frame}

\begin{frame}{Section Summary \& RQ1 Answer}
    \begin{itemize}
        \item CDRF proves cross-modal feasibility; DenoMAE operationalizes self-supervised transfer.
        \item Noise-as-modality + heavy masking deliver low-SNR robustness and vast data efficiency while feeding downstream modules.
        \item \textbf{Answer:} RQ1 is affirmed—vision-based self-supervised learning can approximate world knowledge for RF, yielding a reusable encoder for wireless foundation models.
    \end{itemize}
\end{frame}

